% Part: normal-modal-logic
% Chapter: frame-definability
% Section: definability

\documentclass[../../../include/open-logic-section]{subfiles}

\begin{document}

\olfileid[hi]{nml}{frd}{def}

\olsection{फ़्रेम-परिभाष्यता}

यद्यपि \olref[acc]{thm:soundschemas} के विलोम निहितार्थ विफल होते हैं,
``प्रतिरूप'' को ``फ़्रेम'' से बदलने पर वे मान्य होते हैं:
\olref[acc]{thm:soundschemas} में विचार किए गए गुणों के लिए यह \emph{सत्य}
है कि यदि कोई !!{formula} किसी \emph{फ़्रेम} में मान्य है, तो उस फ़्रेम के
अभिगम्यता संबंध में संगत गुण होता है। अतः विचाराधीन !!{formula}, संगत गुण वाले
फ़्रेमों के वर्गों को \emph{परिभाषित करते हैं}।

\begin{defn}
  यदि $\mClass{F}$ फ़्रेमों का वर्ग है, तो हम कहते हैं कि $!A$,
  $\mClass{F}$ को \emph{परिभाषित करता है}, तभी जब सभी और केवल
  $\mModel{F} \in \mClass{F}$ फ़्रेमों के लिए $\mModel{F} \Entails !A$।
\end{defn}

अब हम फ़्रेमों के लिए पूर्ण परिभाष्यता परिणाम स्थापित करेंगे।

\begin{thm}\ollabel{thm:fullCorrespondence}
यदि \olref[acc]{tab:five} के दाएँ पक्ष का !!{formula} किसी
फ़्रेम~$\mModel{F}$ में मान्य है, तो $\mModel{F}$ में बाएँ पक्ष का गुण है।
\end{thm}

\begin{proof}
  \begin{enumerate}
  \item मान लें \Ax{D}, $\mModel{F} = \tuple{W, R}$ में मान्य है, अर्थात्
    $\mModel{F} \Entails \Box p \lif \Diamond p$। $\mModel{M} =
    \tuple{W, R, V}$ को $\mModel{F}$ पर आधारित प्रतिरूप और $w \in W$ लें।
    हमें ऐसा $v$ दिखाना है कि~$Rwv$। मान लें ऐसा नहीं है: तब प्रत्येक
    $!A$, जिसमें~$p$ भी है, के लिए $\mSat{M}{\Box !A}$ और
    $\mSat/{M}{\Diamond !A}[w]$ दोनों। किंतु तब
    $\mSat/{M}{\Box p \lif \Diamond p}[w]$, जो
    $\mModel{F} \Entails \Box p \lif \Diamond p$ के विरुद्ध है।
  \item मान लें \Ax{T}, $\mModel{F}$ में मान्य है, अर्थात्
    $\mModel{F} \Entails \Box p \lif p$। कोई स्वेच्छ जगत $w \in W$ लें;
    हमें $Rww$ दिखाना है। $u \in V(p)$ तभी रखें जब $Rwu$ (जब $q$, $p$ से
    भिन्न हो, तब $V(q)$ स्वेच्छ है, जैसे $V(q) = \emptyset$)।
    $\mModel{M} = \tuple{W, R, V}$ लें। निर्माण से प्रत्येक $Rwu$ वाले $u$
    के लिए $\mSat{M}{p}[u]$, अतः $\mSat{M}{\Box p}[w]$। पर परिकल्पना से
    $\Box p \lif p$, $w$ पर सत्य है, अतः $\mSat{M}{p}[w]$; और $V$ की
    परिभाषा से यह केवल $Rww$ होने पर संभव है।
  \item हम प्रतिधनात्मक सिद्ध करते हैं। मान लें $\mModel{F}$ सममित नहीं है;
    हम दिखाते हैं कि \Ax{B}, अर्थात् $p \lif \Box\Diamond p$,
    $\mModel{F}= \tuple{W, R}$ में मान्य नहीं है। तब ऐसे $u,v \in W$ हैं
    कि $Ruv$ पर $Rvu$ नहीं। $V$ इस प्रकार परिभाषित करें कि $w \in V(p)$
    तभी जब $Rvw$ नहीं (अन्यथा $V$ स्वेच्छ है)।
    $\mModel{M} =\tuple{W, R,V}$ लें। $V$ की परिभाषा से प्रत्येक ऐसे $w$ के
    लिए $\mSat{M}{p}[w]$ जिसके लिए $Rvw$ नहीं; विशेषतः $Rvu$ न होने से
    $\mSat{M}{p}[u]$। चूँकि $Rvw$ तभी जब $w \notin V(p)$, ऐसा कोई $w$ नहीं
    कि $Rvw$ और $\mSat{M}{p}[w]$; अतः $\mSat/{M}{\Diamond p}[v]$। चूँकि
    $Ruv$, इसलिए $\mSat/{M}{\Box\Diamond p}[u]$ भी। फलतः
    $\mSat/{M}{p \lif \Box\Diamond p}[u]$, अतः \Ax{B}, $\mModel{F}$ में
    मान्य नहीं है।
  \item मान लें \Ax{4}, $\mModel{F} = \tuple{W,R}$ में मान्य है, अर्थात्
    $\mModel{F} \Entails \Box p \lif \Box\Box p$। कोई ऐसे स्वेच्छ जगत
    $u,v,w \in W$ लें कि $Ruv$ और $Rvw$; हमें $Ruw$ दिखाना है। $V$ इस
    प्रकार परिभाषित करें कि $z \in V(p)$ तभी जब $Ruz$ (अन्यथा $V$ स्वेच्छ
    है)। $\mModel{M} =\tuple{W, R, V}$ लें। $V$ की परिभाषा से प्रत्येक
    $Ruz$ वाले $z$ के लिए $\mSat{M}{p}[z]$, अतः $\mSat{M}{\Box p}[u]$।
    पर परिकल्पना से \Ax{4}, $\Box p \lif \Box \Box p$, $u$ पर सत्य है,
    अतः $\mSat{M}{\Box \Box p}[u]$। $Ruv$ और $Rvw$ से
    $\mSat{M}{p}[w]$, और $V$ की परिभाषा से यह केवल $Ruw$ होने पर संभव है।
  \item प्रतिधनात्मक रूप से, मान लें फ़्रेम $\mModel{F} = \tuple{W, R}$
    यूक्लिडीय नहीं है और दिखाएँ कि वह \Ax{5} को असत्य बनाता है, अर्थात्
    $\mModel{F} \Entails/ \Diamond p \lif \Box\Diamond p$। मान लें ऐसे
    जगत $u,v,w \in W$ हैं कि $Rwu$ और $Rwv$ पर $Ruv$ नहीं। प्रत्येक जगत
    $z$ के लिए $V$ इस प्रकार परिभाषित करें कि $z \in V(p)$ तभी जब
    $Ruz$ \emph{नहीं}। $\mModel{M} =\tuple{W, R, V}$ लें। तब परिकल्पना से
    $\mSat{M}{p}[v]$ और $Rwv$ से $\mSat{M}{\Diamond p}[w]$ भी। किंतु ऐसा
    कोई जगत $y$ नहीं कि $Ruy$ और $\mSat{M}{p}[y]$, अतः
    $\mSat/{M}{\Diamond p}[u]$। $Rwu$ से
    $\mSat/{M}{\Box\Diamond p}[w]$, अतः \Ax{5},
    $\Diamond p \lif \Box\Diamond p$, $w$ पर विफल है।
  \end{enumerate}
\end{proof}

\Ax{D} के प्रमाण और अन्य मामलों के बीच अंतर पर ध्यान दें: मूल्यांकन~$V$ का
कोई उल्लेख नहीं हुआ। वस्तुतः हमने सिद्ध किया कि यदि $\mSat{M}{\Ax{D}}$, तो
$\mModel{M}$ सीरियल है। अतः \Ax{D} केवल फ़्रेमों का नहीं, सीरियल
\emph{प्रतिरूपों} का वर्ग परिभाषित करता है।

\begin{cor}\ollabel{prop:D-serial}
  प्रत्येक प्रतिरूप जिसमें \Ax{D} सत्य है, सीरियल है।
\end{cor}

\begin{cor}
\olref[acc]{tab:five} के दाएँ पक्ष का प्रत्येक !!{formula}, बाएँ पक्ष का गुण
रखने वाले फ़्रेमों के वर्ग को परिभाषित करता है।
\end{cor}

\begin{proof}
  \olref[acc]{thm:soundschemas} में हमने सिद्ध किया कि यदि किसी प्रतिरूप में
  बाएँ पक्ष का गुण है, तो दाएँ पक्ष का !!{formula} उसमें सत्य है। अतः यदि किसी
  फ़्रेम~$\mModel{F}$ में बाएँ पक्ष का गुण है, तो दाएँ पक्ष का !!{formula}
  $\mModel{F}$ में मान्य है। \olref{thm:fullCorrespondence} में हमने विलोम
  निहितार्थ सिद्ध किए: यदि दाएँ पक्ष का कोई !!{formula} $\mModel{F}$ में मान्य
  है, तो $\mModel{F}$ में बाएँ पक्ष का गुण है।
\end{proof}

\begin{prob}
दिखाइए कि यदि \olref[nml][frd][acc]{tab:anotherfive} के दाएँ पक्ष का
!!{formula} किसी फ़्रेम~$\mModel{F}$ में मान्य है, तो $\mModel{F}$ में बाएँ
पक्ष का गुण है। इसके लिए ऐसे फ़्रेम पर विचार करें जो बाएँ पक्ष का गुण
\emph{संतुष्ट नहीं करता}, और उपयुक्त~$V$ परिभाषित करें जिससे दाएँ पक्ष का
!!{formula} किसी जगत पर असत्य हो।
\end{prob}

\olref{thm:fullCorrespondence} यह भी दिखाता है कि गुणों को संयोजित किया जा
सकता है: उदाहरणार्थ, यदि \Ax{B} और \Ax{4} दोनों $\mModel{F}$ में मान्य हैं,
तो फ़्रेम सममित और संक्रमणशील दोनों है। अनेक महत्त्वपूर्ण मोडल तर्कशास्त्र
कुछ फ़्रेम-गुणों को संयोजित करने वाले सभी फ़्रेमों में मान्य !!{formula}ों के
समुच्चय के रूप में अभिलक्षित होते हैं। अतः हम उन्हें उन सभी फ़्रेमों में
मान्य !!{formula}ों के समुच्चय के रूप में अभिलक्षित कर सकते हैं जिनमें संगत
परिभाषक !!{formula} मान्य हैं। उदाहरणार्थ, शास्त्रीय तंत्र \Log{S4}, सभी
स्वपरावर्ती और संक्रमणशील फ़्रेमों में मान्य !!{formula}ों का समुच्चय है,
अर्थात् उन सभी में जहाँ \Ax{T} और~\Ax{4} दोनों मान्य हैं। \Log{S5}, सभी
स्वपरावर्ती, सममित और यूक्लिडीय फ़्रेमों में मान्य सूत्रों का समुच्चय है,
अर्थात् उन सभी में जहाँ \Ax{T}, \Ax{B} और \Ax{5} सभी मान्य हैं।

$R$ के गुणों के बीच तार्किक संबंध सामान्यतः उनके संगत परिभाषक !!{formula}ों
के बीच संबंधों के अनुरूप होते हैं। उदाहरणार्थ, प्रत्येक स्वपरावर्ती संबंध
सीरियल है; अतः जब भी \Ax{T} किसी फ़्रेम में मान्य है, \Ax{D} भी मान्य है।
(ध्यान दें कि यह संबंध परिणाम का संबंध \emph{नहीं} है। ऐसा नहीं कि जब भी
$\mSat{M}{\Ax{T}}[w]$, तब $\mSat{M}{\Ax{D}}[w]$।) हम ऐसे कुछ संबंध दर्ज
करते हैं।

\begin{prop}\ollabel{prop:relation-facts}
  मान लें $R$ किसी समुच्चय $W$ पर द्विपदी संबंध है; तब:
  \begin{enumerate}
  \item यदि $R$ स्वपरावर्ती है, तो वह सीरियल है।
  \item यदि $R$ सममित है, तो वह संक्रमणशील तभी है जब वह यूक्लिडीय है।
  \item यदि $R$ सममित या यूक्लिडीय है, तो वह दुर्बलतः दिष्ट है (उसमें
    ``हीरक गुण'' है)।
  \item यदि $R$ यूक्लिडीय है, तो वह दुर्बलतः संबद्ध है।
  \item यदि $R$ फलनीय है, तो वह सीरियल है।
  \end{enumerate}
\end{prop}

\begin{prob}
  \olref[nml][frd][def]{prop:relation-facts} सिद्ध कीजिए।
\end{prob}

\end{document}
